%% kapa: exam=1 round=26 year=2015 period=1 form=A subject=civil_law subject_ko=민법 seq=27 number=27
%% source: Q-Net 감정평가사 기출문제 (cst00309) artlSeq=5138001 file=extracted/1008/5138001/제26회_감정평가사 1차_기출/제26회_감정평가사 1차_기출/1교시_감평_A형.pdf
%% conversion: kordoc 4.24.51 -> tools/gen1.py (auto); review=proofread-sonnet; flags=none
\begin{question}{27}
혼동으로 물권이 소멸하는 경우는? (다툼이 있으면 판례에 따름)
\begin{choices}
  \item 甲의 토지에 乙이 1번 저당권, 丙이 2번 저당권을 취득한 후 乙이 토지 소유권을 취득하는 경우
  \item 甲의 건물에 乙이 저당권을 취득한 다음 그 건물을 매수하여 소유권이전등기를 마쳤는데, 그 매매계약이 원인무효임이 밝혀진 경우
  \item 甲의 건물에 乙이 1번 저당권, 丙이 2번 저당권을 취득한 후 丙이 건물 소유권을 취득하는 경우
  \item 甲의 토지에 乙이 지상권을 취득하고, 그 지상권 위에 丙이 저당권을 취득한 후 乙이 토지 소유권을 취득하는 경우
  \item 甲의 토지에 대한 乙의 지상권 위에 丙이 1번 저당권, 丁이 2번 저당권을 취득한 뒤 丙이 乙의 지상권을 취득하는 경우
\end{choices}
\end{question}
\answer{3}
